\documentclass[12pt]{article}
\input{.house-style/preamble.tex}
\usepackage{forest}
\usepackage{xurl}
\usepackage{float}
\useforestlibrary{linguistics}
\setlength{\headheight}{14pt}
\clubpenalty=10000
\widowpenalty=10000
\AtBeginBibliography{\emergencystretch=1em}
\setcounter{biburlnumpenalty}{100}
% Local display normalization; the shared bibliography remains the source.
\DeclareSourcemap{\maps[datatype=bibtex]{\map{
  \step[fieldsource=journal, match=\regexp{^Glossa:\s+a\s+journal\s+of\s+general\s+linguistics$},
        replace={Glossa: A Journal of General Linguistics}]
}}}
\newcommand{\synnode}[2]{\shortstack{\scriptsize #1\\#2}}
\forestset{head edge/.style={edge={line width=1pt}}, nominal tree/.style={for tree={align=center, parent anchor=south, child anchor=north, l sep=5mm, s sep=4mm, inner sep=1.5pt}}}
\hypersetup{pdftitle={Determinatives as nouns in English},pdfkeywords={determinatives, nouns, lexical categories, noun phrases, English}}
\title{Determinatives as nouns in English}
\author{Brett Reynolds \orcidlink{0000-0003-0073-7195}%
\thanks{Contact: \href{mailto:brett.reynolds@humber.ca}{brett.reynolds@humber.ca}}\\
Humber Polytechnic \& University of Toronto}
\date{Draft, September 2026}
\begin{document}
\maketitle

\begin{abstract}
I argue that English determinatives, including articles, demonstratives, and quantifiers, belong within Noun alongside common nouns, proper nouns, and pronouns. Quantificational common nouns provide the closest comparison: they share complement patterns, number transparency, and restricted dependents with parts of the determinative inventory. The strongest adjectival counterweight connects grade, degree modification, and comparative complementation in four quantifiers. Number and referential contrasts supply further nominal connections. Restricted articles belong to the grouping through their integration into the determinative system. In \mention{take some apples} and \mention{take some}, \mention{some} belongs to the same category. For the independent use, I favour ordinary Head relations within the NP, rather than \textit{CGEL}'s fusion of determiner and Head functions on one constituent. This gives determinatives the ordinary noun-phrase structure that their inclusion within Noun makes natural. Neither choice requires the other, but their fit supports the package.
\end{abstract}

\noindent\textbf{Keywords:} determinatives; nouns; lexical categories; noun phrases; English

\par\medskip
\noindent{\small\textsc{AI use.} For the September 2026 revision, GPT-6 (Astra), Claude Fable~5.1, Claude Opus~5, Claude Haiku~4.5, and GLM-5.3-Flash assisted drafting, source retrieval, script development, or critical review; GPT-5.6 (Sol) checked numerical outputs. The accompanying materials record the tools and their uses. Responsibility for the analysis is mine.\par}
\medskip

\section{The question and the alternatives}\label{sec:intro}

What's the categorial relationship among words such as \mention{some}, \mention{me}, \mention{apple}, and \mention{Brett}? I argue that all four are nouns: determinatives form a coordinate subcategory alongside common nouns, proper nouns, and pronouns. I write the superordinate category as Noun.

I adopt the general framework of \textit{The Cambridge grammar of the English language} (\textit{CGEL}; \citealp{huddleston2002}). Unlike \textit{CGEL}, though, I include determinatives within Noun. The claims concern synchronic English lexical categories.\footnote{I treat these categories as language-specific: a categorization supported by another language's grammar doesn't determine the English categorization.}

I use \term{determinative} for the category containing articles, demonstratives, and quantifiers such as \mention{the}, \mention{this}, \mention{some}, \mention{every}, and \mention{many}.\footnote{For a fuller inventory, see the online \href{https://www.cambridge.org/highereducation/api/resources/2AC8DF5CA558F8A6FCDFACD0EAF8EA08}{\enquote{List of determinatives in English}} accompanying \textcite{Huddleston2021}.} I reserve \term{determiner} for the syntactic function within the noun phrase (NP) characteristically performed by phrases headed by these words. This distinction separates what kind of word \mention{some} is from what syntactic relationships its phrase participates in.

Compare \mention{take some apples} with \mention{take some}. The word \mention{some} remains determinative in both. Before \mention{apples}, its phrase functions as determiner; without \mention{apples}, the whole expression functions as object.

In these NP comparisons, \term{dependent} describes a determinative's use as determiner or internal modifier with a nominal target; \term{independent} describes use without that target. An independent expression can still have dependents of its own, as in \mention{some of the wine}, and rely on context for its interpretation.

Two choices are involved. The first is whether determinatives belong within Noun. The second concerns the Head relations within independent NPs. In \textit{CGEL}, the determinative phrase (DP) headed by independent \mention{some} jointly fills determiner and Head, a \term{fusion of functions} \citep[410--412]{huddleston2002}. The alternative has \mention{some} head a nominal (Nom) directly, with Nom heading NP.

A standard tree has unique parentage: each constituent below the root has one parent. Fusion allows branches to converge on a shared constituent with two parents. In the \term{ordinary-Head analysis}, the determinative's projection needs no such join. Figure~\ref{fig:some} in §\ref{sec:head-relations} compares the structures.\footnote{The proposal retains fusion elsewhere, including Mod--Head fusion in \mention{the rich}. The preference for unique parentage concerns the determinative's Head configuration, rather than every construction in the grammar.}

I call the inclusion of determinatives within Noun the \term{D-noun categorization}. The \term{D-noun analysis} proposed here combines it with ordinary headedness. Noun membership makes ordinary nominal structure a natural treatment; that structural regularity strengthens the case for the grouping.

Neither commitment requires the other: a noun's projection can fill a fused function, and a grammar can keep determinative (D) separate while assigning ordinary Head. Table~\ref{tab:accounts} shows the four combinations.\footnote{Another logical alternative posits an unrealized noun as Head. I set that option aside here. \textit{CGEL} discusses the limits of ellipsis as a general account \citep[420--421]{huddleston2002}.}

\begin{table}[H]
\centering\small
\caption{Four combinations of categorization and Head analysis. The proposed package is the first row. All keep \mention{apples} as the ultimate lexical head of \mention{some apples}.}\label{tab:accounts}
\begin{tabular}{>{\raggedright\arraybackslash}p{3.6cm}>{\raggedright\arraybackslash}p{4.1cm}>{\raggedright\arraybackslash}p{4.4cm}}
\toprule
Account & Place of determinative & Independent \mention{some} \\
\midrule
D-noun, ordinary Head & A subcategory of Noun & Ordinary Head through nominal projection \\
Separate D, ordinary Head & A primary category alongside Noun & Ordinary Head, licensed for both Noun and D \\
Separate D, fused Head & A primary category alongside Noun & Joint determiner and Head functions \\
D-noun, fused Head & A subcategory of Noun & Joint determiner and Head functions \\
\bottomrule
\end{tabular}
\end{table}

Both Head analyses give the whole independent expression NP status. Their difference concerns its internal structure. In the separate-D ordinary-Head implementation, words of category D project NP (§\ref{sec:det-uniform}).

Lexical categories and syntactic functions cut across one another. An NP can function as determiner, as in \mention{\underline{Kim's} book}, and a determinative-headed phrase can function as modifier, as in \mention{the \underline{many} people} \citep{payne2010,pullummiller2022nps}.

The DP hypothesis of \textcite{abney1987} concerns a different Head relation. Under that hypothesis, D heads expressions such as \mention{some apples}. I retain noun-headed NPs and reject the DP analysis for English, following the arguments of \textcite{pullummiller2022nps} and \textcite{bruening2020nominal}. Here DP means \term{determinative phrase}, as in \textit{CGEL}.

Category membership doesn't remove lexical restrictions, such as the one applying to \mention{every}. \mention{Every apple} is grammatical, but unlike \mention{some} in \mention{I'll take some}, \mention{every} can't occur independently in \ungram{\mention{I'll take every}}. The exclusion rests on absence where opportunity is frequent: independent \mention{each} is common in the same corpora, and across roughly a thousand COCA and NOW lines with \mention{every} before a verb it never stands as an ordinary fused head. Where it does stand without a noun, the noun is recoverable from a coordinate or the preceding turn, or \mention{every} closes a binomial.\footnote{COCA: \enquote{if not every, then almost every male head of household} (\textit{Social History}, 1995); \enquote{But every man will want them. \ldots{} Not every.} (fiction, 2002); \enquote{to be open and generous to all and every} (fiction, 2015); \mention{any and every} (blog, 2012). Of 618 COCA lines for \mention{every} before a verb and 164 NOW lines for \mention{every was} and \mention{every has}, none shows the fused-head use; informal text supplies \mention{every} as a slip for \mention{everyone}. The screened lines are under \texttt{corpus/exclusion\_tests/}, and the check is tabled in the quantifier supplement.}

\textit{CGEL} already includes pronouns within Noun on the basis of their phrases' functions, despite differences from common and proper nouns in inflection and dependents \citep[327--328]{huddleston2002}. The inclusion of auxiliaries within Verb supplies a further precedent for preserving distinctive properties within a broader lexical category \citep{pullumwilson1977}. These precedents motivate comparing the whole grammatical profile, rather than requiring every member to display each nominal property.

The comparison weighs three considerations: how broadly a property occurs, how it recurs across constructions, and how specifically it connects the groups. These considerations apply equally to nominal and adjectival connections. A pattern concentrated in a few words can be substantial evidence; its contribution depends on its grammatical relationships as well as its reach.

Earlier accounts connect articles with pronouns or give both nominal structure, at different analytical levels (Appendix~\ref{sec:historical}). Hudson's nested noun taxonomy is considered in §\ref{sec:inheritance}; Spinillo's restricted article grouping in §\ref{sec:articles}.

Section~\ref{sec:proposal} compares the grammatical profiles, and §\ref{sec:articles} completes the case for the intended inventory. Section~\ref{sec:evidence} then compares internal structures, which the fragment in §\ref{sec:economy} makes explicit and assesses. The further question of coordinate or nested rank follows in §\ref{sec:inheritance}.

\section{Grounds for a broader Noun category}\label{sec:proposal}

The existing Noun category contains common nouns, proper nouns, and pronouns with markedly different meanings, inflection, and dependents. Ordinary count nouns alone are an inadequate comparison group. The question is how determinatives connect with that broader range, and how those connections compare with their affinities to adjectives.

\subsection{Functions and constructional range}\label{sec:independent}

The external syntax of a phrase concerns the functions it fills in larger constructions. Noun phrases can serve as subjects, objects, and complements of prepositions. Independent determinative expressions enter these constructions too, with the lexical restrictions illustrated below.

Independent use is widespread within the determinative inventory. \textit{CGEL} discusses such uses for demonstratives and quantifiers including \mention{some}, \mention{all}, \mention{both}, \mention{many}, \mention{few}, \mention{several}, \mention{each}, \mention{either}, \mention{neither}, \mention{much}, and \mention{enough}, while recording lexical restrictions and the separate forms \mention{no}/\mention{none} \citep[371--372, 410--424]{huddleston2002}. The generalization concerns the availability of independent constructions across a lexical category, not unrestricted acceptability in every sentence frame.\footnote{Material before a determiner falls outside the independent-use comparison. \textit{CGEL} treats \mention{all}/\mention{both} in \mention{all/both the books} as predeterminer modifiers, \mention{quite}/\mention{rather} before \mention{a good idea} as peripheral modifiers, and \mention{such}/exclamative \mention{what} before \mention{a disaster} as adjectives \citep[433--437]{huddleston2002}. The fixed \mention{many a} is a complex determinative restricted to Det function \citep[394]{huddleston2002}. The \mention{half} in \mention{half a cake} is a common noun used as a predeterminer modifier \citep[434]{huddleston2002}. These constructions don't add evidence for independent determinative heads.}

In the constructed examples in (\ref{ex:external}), compare the bracketed NPs as subjects, objects, and complements of prepositions. Assume that a woman named Kim and a group of people are already under discussion.

\begin{samepage}

\ea\label{ex:external}
\ea \mention{\textup{[}People\textup{]} left.}\qquad \mention{I see \textup{[}people\textup{]}.}\qquad \mention{with \textup{[}people\textup{]}}
\ex \mention{\textup{[}Kim\textup{]} left.}\qquad \mention{I see \textup{[}Kim\textup{]}.}\qquad \mention{with \textup{[}Kim\textup{]}}
\ex \mention{\textup{[}She\textup{]} left.}\qquad \mention{I see \textup{[}her\textup{]}.}\qquad \mention{with \textup{[}her\textup{]}}
\ex \mention{\textup{[}Some\textup{]} left.}\qquad \mention{I see \textup{[}some\textup{]}.}\qquad \mention{with \textup{[}some\textup{]}}
\z\z

\end{samepage}

Independent uses containing fused-head adjective phrases (AdjPs) prevent a simple inference from these positions to nounhood. \mention{The rich} and \mention{the poor} can fill nominal argument positions. Comparative and superlative adjectives also head expressions without a generic human interpretation: \textit{CGEL}'s \mention{the most important of her criticisms} is an NP containing a partitive \mention{of}-phrase \citep[332--333, 416--423]{huddleston2002}.

Independent \mention{some} can form a one-word NP, whereas an NP with \mention{rich} as fused head requires the definite article \mention{the} on the generic human reading \citep[417--418]{huddleston2002}. That's a local contrast: argument NPs headed by singular count common nouns also need determination.

Existentials add a construction to this comparison. With possible solutions under discussion, the constructed \mention{There are several} places the independent quantifier in displaced-subject position; dummy \mention{there} is the subject \citep[1391--1393]{huddleston2002}. Section~\ref{sec:quant-controls} compares this use with quantity words near the adjective--determinative boundary, distinguishing expressions with partitive or relative dependents from those interpreted through context alone.

Determiner function also cuts across the existing noun subcategories. In \mention{my preferences}, \mention{Kim's preferences}, and \mention{people's preferences}, the determiner is an NP ultimately headed by a pronoun, a proper noun, and a common noun respectively \citep[354--355, 470--471]{huddleston2002}.\footnote{\textit{CGEL} assigns these genitives the combined function Subject--Det \citep[472--473]{huddleston2002}. I treat them as Det here, without the additional subject function.}

The compound determinative \mention{someone}, following the categorization of \textcite[§1.3]{Payne2007}, participates in \mention{someone's preferences}. Plain-case NPs also fill Det: \mention{what size} in \mention{what size shoes}, \mention{that size} in \mention{that size shoes}, and \mention{Sunday} in \mention{Sunday morning} \citep[356]{huddleston2002}.

Determiner remains the characteristic function of determinative phrases. This specialization could support a separate primary category. Its tasks are themselves closely connected with nominal reference: Det marks definiteness and often contributes quantification \citep[354--359]{huddleston2002}. The issue is whether that specialization warrants a primary boundary or a distinction within Noun.

The existing NP determiners provide a positive comparison. In \mention{Kim's preferences}, the embedded NP supplies an identifying anchor through its own referent; \mention{Sunday} and \mention{what size} specify a day or a dimension. Determinatives' deictic and quantitative specifications fit this nominal pattern. I take that fit to support treating their specialization in Det as a distinction within Noun.

Preposition phrases (PPs) remain a restricted alternative, as in \mention{up to twenty minutes} and \mention{between fifty and sixty tanks} \citep[356]{huddleston2002}. Number, countability, and other selectional conditions distinguish the determining expressions.

Modifier and adjunct functions also cut across the four noun groups. Compare the modifiers in \mention{dog houses}, \mention{Canada Day}, \mention{the manager herself}, and \mention{the few people}. Emphatic \mention{herself} also functions as a clause adjunct in \mention{The manager detected the error herself} \citep[1496--1497]{huddleston2002}. Temporal NPs such as \mention{that day} and \mention{Sunday} supply adjuncts, as does degree \mention{enough} in \mention{I hadn't prepared enough} (§\ref{sec:enough}). These are shared functions with construction-specific distributions.

Relative-clause postmodification supplies a further constructional comparison. Common nouns freely take integrated relatives, as in \mention{people who came}; personal pronouns permit a restricted range, including \mention{we who have read the report} \citep[430]{huddleston2002}. Determinative examples include \mention{few who come ever leave}, \mention{those who came}, \mention{that which remains}, and \mention{something that you need to know}.

Compounds also permit \mention{anyone who asks} and \mention{everything that matters}. With books under discussion, compare \mention{some that I saw} and \mention{two that I have seen}. \textit{CGEL} describes relative postmodification with demonstratives and compounds \citep[414, 422--423]{huddleston2002}.

Compound determinatives also take post-head adjectives, as in \mention{I need something reliable and good looking}. The position and interpretation of these modifiers have specialized conditions (§\ref{sec:compounds}).

\subsection{The connection from quantificational common nouns}\label{sec:quant-nouns}

Quantificational common nouns make the connection more specific than shared quantity meanings. I use \term{quantifier} across lexical categories for expressions specifying quantity. In \mention{a lot of the delegates} and \mention{many of the delegates}, both heads quantify over a partitive \term{domain}: the whole from which a quantity is drawn, here the delegates. But \mention{lot} also permits \mention{a lot of delegates}, whereas \ungram{\mention{many of delegates}} is excluded \citep[349]{huddleston2002}.

Agreement provides another connection. Compare the constructed \mention{A lot of the people were waiting} and \mention{Some of the people were waiting}, with plural agreement, against \mention{A lot of the work was finished} and \mention{Some of the work was finished}, with singular agreement. The whole subject NP's number depends on the complement of \mention{of}. \textit{CGEL} calls quantificational \mention{lot} \term{number-transparent} \citep[349--350, 411--412]{huddleston2002}.

\mention{Each of the people} takes singular agreement by default, and ordinary \mention{a photograph of the people} takes it categorically, so neither is number-transparent in \textit{CGEL}'s sense. The contrast isn't absolute, though. Plural agreement after \mention{each of them} occurs in about a sixth of COCA tokens, and after \mention{each of the} with a plural noun in about a third, in edited prose as well as speech.\footnote{COCA List counts (\href{https://www.english-corpora.org/coca/}{english-corpora.org/coca}, 17 September 2026): \mention{each of them} + \mention{is/has/was} 389, + \mention{have/are/were} 78; \mention{each of the} + plural noun + \mention{was} 151, + \mention{were} 73 over 173 noun-specific strings, a share inflated by sentences in which the verb agrees with another head. Attestations: \enquote{Each of them have had a career that is the stuff of dreams} (\textit{USA Today}, 2018); \enquote{Each of the participants were administered the survey packet} (\textit{Professional School Counseling}, 2007); \enquote{Each of the films were selected from a collection of family films} (\textit{Roeper Review}, 2002). Queries and screened lines are in the project repository under \texttt{corpus/exclusion\_tests/}.} That variation is agreement with the nearer or notional plural, the pattern familiar from \mention{a number of the people were}, and it leaves \mention{each} singular; with \mention{some of the people}, by contrast, plural agreement is the only option. Number transparency specifically connects quantificational \mention{lot} with number-neutral \mention{some}.

Restricted dependents also occur on the common-noun side. \textit{CGEL} categorizes quantificational \mention{plenty} as a common noun whose use resists determination and modification; \mention{lot} requires \mention{a} and permits only limited modification \citep[349--350]{huddleston2002}. Established common nouns thus approach the determinative profile as determinatives approach theirs.

The connection also extends beyond argument NPs. Quantificational nouns occur in degree modifiers such as \mention{a great deal smaller} and \mention{plenty big enough} \citep[549--550]{huddleston2002}. Section~\ref{sec:enough} compares these with determinative degree modifiers, including their attributive distribution. Meaning, complementation, agreement, and restricted dependents give the proposed grouping a specific basis within established Noun.

Common nouns also permit a wider range of PP and clausal complements, whereas pronouns and primary naming uses of proper nouns have much more restricted dependents \citep[429--430, 439--443, 517--521]{huddleston2002}. The shared partitive in \mention{a lot/some of the wine} connects these different complement systems.

\subsection{Quantificational adjectives and nominal independence}\label{sec:quant-controls}

\mention{Numerous people} and \mention{many people} both quantify, but their syntax differs. Degree modification exposes a contrast: \mention{so many mistakes} is possible, while \ungram{\mention{so numerous mistakes}} is excluded \citep[539--540]{huddleston2002}. Quantity meaning and prenominal position provide a semantic control for the nominal comparison. Independent use tests whether their further distributions sustain the distinction.

The boundary already cuts through this semantic range in \textit{CGEL}. It includes \mention{several} among determinatives and treats quantificational \mention{certain} and \mention{various} as marginal members, partly on the evidence of partitives. Their \mention{of}-phrases are hardly omissible, with speaker and register restrictions on the partitive uses \citep[392--393, 411--413]{huddleston2002}. Predicative \mention{Its advantages are several} and \mention{Their enemies were many} provide further overlap, though the latter is uncommon and formal \citep[392, 395--396]{huddleston2002}.

Independent uses extend to \mention{numerous}, \mention{multiple}, and \mention{countless}. Attested examples include subject \mention{numerous were injured}, anaphoric \mention{are there multiple?}, and \mention{paving the way for countless who followed him}, where a relative postmodifies the complement of \mention{for}.\footnote{The accompanying \href{run:quantifier-controls.pdf}{\textit{Quantificational controls: attestations and provenance}} preserves the comparison table, attestations, source locations, and verification limits. Descriptions in \textit{CGEL}, constructed comparisons, positive attestations, and searches with no qualifying result are distinguished. A second table there records corpus checks of the exclusions this paper states, with the search strings, hit counts, and outcomes.} These selected occurrences establish overlap, without establishing equal acceptability or frequency. Assigning every independent quantity word to determinative would make independent use a circular category diagnostic.

Existential occurrence doesn't decide the structural analysis either. Adjective fusion permits indefinite ordinal \mention{a second} \citep[416]{huddleston2002}, giving the constructed comparison \mention{There was a second}. Plural agreement likewise occurs with adjective-containing \mention{the rich}, whose adjective lacks number inflection \citep[418]{huddleston2002}. The number transparency of \mention{some} and \mention{a lot} in §\ref{sec:quant-nouns} supplies a more specific comparison, with explicit count and non-count domains.

The categorial significance of independent use depends on its combination with degree syntax, determination, and dependents.

\subsection{The connected adjectival profile}\label{sec:adjectival-profile}

The strongest adjectival connection also concerns a restricted group: \mention{few}, \mention{many}, \mention{much}, and \mention{little}. They distinguish plain, comparative, and superlative forms: \mention{few}/\allowbreak\mention{fewer}/\allowbreak\mention{fewest}, \mention{many}/\allowbreak\mention{more}/\allowbreak\mention{most}, and corresponding paradigms for \mention{little} and \mention{much} \citep[391--395]{huddleston2002}. Grade combines with degree-modifier selection and comparative complementation, as in \mention{more than ten}. The three properties form a connected profile.

Degree \mention{very}, \mention{so}, and \mention{too} select gradable heads: \mention{very tall} and \mention{very few}, but \ungram{\mention{so you}}, \ungram{\mention{so the book}}, and \ungram{\mention{too my life}}. These degree modifiers also exclude non-gradable determinatives: \ungram{\mention{very every}}, \ungram{\mention{very some}}, and \ungram{\mention{very this}}. The associated series includes \mention{so few}, \mention{too few}, \mention{how few}, and \mention{as few as}. These degree patterns select \mention{few}, \mention{many}, \mention{much}, and \mention{little}, the four determinatives that inflect for grade \citep[393--395, 431--432]{huddleston2002}.\footnote{The restriction concerns this degree series. \textit{CGEL} also records modifiers of completeness, including \mention{marginally enough} and \mention{absolutely all}, with different selectional ranges \citep[432]{huddleston2002}. Those combinations don't extend the \mention{very}/\mention{so}/\mention{too}/\mention{how} series beyond the four degree determinatives.}

With the complex determinative \mention{a few}, modifier position distinguishes two kinds of attachment. \textit{CGEL} explicitly contrasts internal \mention{very} in \mention{a very few mistakes} with peripheral \mention{quite} in \mention{quite a few mistakes} \citep[p.~392, (61)]{huddleston2002}. The reversed orders \ungram{\mention{very a few}} and \ungram{\mention{a quite few}} are excluded in these constructions. Here \mention{a} belongs to the complex determinative; \mention{very} occurs within it and \mention{quite} outside it.

Approximative and focusing modifiers have a broader distribution. Compare \mention{almost every}, \mention{nearly all}, \mention{hardly any}, \mention{practically no}, \mention{almost everyone}, and \mention{almost ten}. They have counterparts in established NP modification: \mention{almost the whole class}, \mention{hardly a day}, and \mention{only the best}. In \mention{almost every student}, \mention{almost} approximates the quantity; it doesn't intensify a gradable property of \mention{every}.

The broader approximative pattern supplies NP parallels distinct from the four quantifiers' adjectival profile. Section~\ref{sec:payne} develops an attachment analysis of that contrast, available under either ordinary-Head taxonomy.

Overlap with adjectives likewise needs phrase levels kept distinct. Both \mention{the people} and \mention{the rich} are NPs in \textit{CGEL}, though only the former has a noun as lexical head. In \mention{a soccer/round ball}, a noun-headed expression and an AdjP share modifier function; in \mention{Jones became president/ill}, NP and AdjP share predicative complement function. External functions alone therefore don't establish nounhood.

\mention{She is a beauty} ascribes a property, whereas \mention{That is Kim} identifies a person. \textit{CGEL} excludes specifying \mention{be} clauses from its noun--adjective diagnostic because phrases from other categories occur in them too \citep[536]{huddleston2002}.

NP grammar also distinguishes predicative and argument uses. A bare-role NP such as \mention{president} is licensed in \mention{I'd like to be president}, but requires determination in the corresponding object use \mention{I'd like to meet the president} \citep[328]{huddleston2002}. Determinatives' asymmetries fit this broader variation within Noun.

Bare colour expressions and evaluative comparative subjects complicate a simple distributional boundary. \textit{CGEL}'s \mention{Henrietta likes red shirts, and I like blue} permits reduction licensed by coordination. The constructed \mention{Blue is good} raises a separate ambiguity between noun and adjective, while the attested \mention{Bluer is better} establishes a bare comparative subject without settling its phrase category.\footnote{The \mention{blue}, \mention{old}, and \mention{small} examples on \textit{CGEL} p.~417 all occur in coordinated contrasts. \mention{Bluer is better} accompanies a results colour scale in Stanford CS329X, \href{https://web.stanford.edu/class/cs329x/slides/Lecture12_B_Codeswitching\%20LLMs.pdf}{\enquote{Codeswitching LLMs}}, slide~8 (accessed 10 September 2026). I interpret it as evaluating a degree of blueness. Subject function alone doesn't establish NP status \citep[236]{huddleston2002}.} These cases leave the wider comparison of argument distribution and head properties necessary.

\subsection{Number, genitive marking, and reference}\label{sec:inflection}

Inflection supplies several connections with established nouns, though no one contrast runs through the whole category. Common nouns typically distinguish singular and plural; proper nouns permit number inflection in restricted uses such as \mention{the Smiths}. Demonstratives likewise distinguish \mention{this}/\mention{these} and \mention{that}/\mention{those} \citep[373, 521]{huddleston2002}. These are contrasts between forms of a lexeme, unlike the fixed number restrictions of \mention{each} and \mention{several}.

Number needn't be inflectional wherever it matters. \textit{CGEL} treats singular and plural \mention{you} as distinct lexemes, differentiated overtly only by the reflexive forms \mention{yourself} and \mention{yourselves} \citep[486]{huddleston2002}. Number-neutral \mention{some} has no corresponding singular--plural alternation. The demonstratives supply positive morphological evidence; the absence of that alternation elsewhere doesn't distinguish determinatives from pronouns as whole categories.

The demonstrative number contrasts remain visible in independent uses. Generic human \mention{the rich}, by contrast, is a plural NP whose adjective lacks number inflection and still takes adverb modifiers: \mention{the very rich} \citep[418]{huddleston2002}. The difference concerns overt paradigms, rather than the mere possession of a number value by the whole NP.

Case gives a further partial connection. Genitives occur across the proposed subcategories: \mention{dog's}, \mention{Kim's}, \mention{my}, and \mention{someone's}. Personal pronouns have fuller case paradigms. Most determinatives lack case inflection, and compounds such as \mention{someone} and \mention{something} may owe their genitives to the nominal component \citep[423--424, 479--480]{huddleston2002}. This supports the compounds' nominal affinity more directly than the category's membership as a whole.

Meaning and reference vary within Noun. \mention{Apple} supplies descriptive content; \mention{Kim}, in its primary naming use, identifies through a name; \mention{she} depends on context. Many determinatives also have pro-form uses. They range from deictic \mention{this} to quantitative \mention{many} and universal \mention{every} \citep[358--360, 370--405, 425--428, 515--521]{huddleston2002}. Deixis connects demonstratives with pronouns, while quantity connects determinatives with common nouns such as \mention{number} and \mention{majority}.

These meanings aren't exclusive to nouns. Adjectives such as \mention{singular} and \mention{plural} also concern number; \mention{proximate} and \mention{distal} describe spatial relations relevant to demonstrative contrasts, though spatial meaning needn't itself be deictic. Nor does Noun imply close semantic similarity throughout: \mention{every} quantifies over a nominal restriction, while \mention{Kim} identifies an individual. Semantic affinities contribute to the comparison without settling its boundaries.

Pro-form gender \citep{reynolds2025proformgender} connects these ways of referring. Descriptions such as \mention{the woman} and names such as \mention{Kim} identify referents whose construal bears on pro-form choice. Pronouns and determinatives express contrasts between personal \mention{she}/\mention{somebody} and non-personal \mention{it}/\mention{something}. Relative \mention{who}/\mention{which} supplies another contrast.\footnote{Whether independent \mention{what} and relative \mention{which} belong to pronoun or determinative leaves them within the proposed Noun category. I leave that internal boundary open here; \textit{CGEL}'s inventory suffices for the present comparison \citep[397--399]{huddleston2002}.}

Interrogative and relative properties are distinct from functions such as object. An interrogative object may be headed by a pronoun (\mention{who}) or a common noun (\mention{which book}); the latter obtains its interrogative property from a dependent. Neither the object function nor the interrogative construction determines the lexical head's category.

Inflection and reference provide partial connections, with different reach. Demonstrative number directly compares forms of determinative lexemes with nominal number paradigms; compound genitives may depend on their nominal component. The \mention{no}/\allowbreak\mention{none} alternation concerns form selection, discussed with restricted membership in §\ref{sec:form-selection}.\footnote{Derivations such as \mention{nothingness} and \mention{oneness} are weaker category diagnostics: \mention{-ness} also attaches to adjectives and other bases \citep[Ch.~19, §5.7.2(i)]{huddleston2002}. Numeral morphology likewise needs its inputs distinguished. \textcite[§§5.4--5.5]{reynolds2026numerals} derives fractional nouns from cardinal nouns and treats \mention{two thousand and twenty-seventh} as a coordination whose final coordinate is adjectival. Neither process independently establishes determinative nounhood.}

\subsection{Weighing the profiles}\label{sec:membership}

Table~\ref{tab:existing} summarizes the range within each category; individual lexemes differ in which properties they exhibit. Functions concern phrases, inflection concerns word forms, and inventory closure concerns categories. AdjP and AdvP denote adjective and adverb phrases. The modification rows identify proposed determinative attachments, explained in §\ref{sec:payne}. The rows aren't equally weighted tests, and several properties belong to one connected pattern.

Proper names permit nominal premodifiers such as \mention{architect Norman Foster}; \mention{all} in \mention{all we who signed up} is peripheral, rather than an internal pronoun modifier \citep[429--430, 519--520]{huddleston2002}. The cardinal examples are subject to the category qualifications in §§\ref{sec:evidence} and~\ref{sec:costs}.

\begin{table}[H]
\centering\small
\caption{Profiles compared. Determinative modifier attachments are analytical commitments; interrogative and relative forms follow \textit{CGEL}'s inventory.}\label{tab:existing}
\setlength{\tabcolsep}{4pt}
\begin{tabular}{>{\raggedright\arraybackslash}p{2.25cm}*{4}{>{\raggedright\arraybackslash}p{2.85cm}}}
\toprule
Dimension & Common noun & Proper noun & Pronoun & Determinative \\
\midrule
Argument / Det functions & NP arguments; genitive and restricted plain Det & NP arguments; genitive and restricted plain Det & NP arguments; genitive Det & NP arguments; plain and genitive Det \\
Accepts determination & Broad contrasts; singular count arguments normally require it & Restricted in primary naming uses & Normally excluded & Lexically restricted: \mention{the few}, \mention{these three} \\
Internal modification & Productive AdjP and nominal premodifiers, including determinative-headed phrases; relatives & Restricted AdjP and nominal premodifiers; embellishments & Restricted AdjP premodifiers and relative postmodifiers & Degree AdvP on four quantifiers; approximatives on cardinals; NP and determinative-headed modifiers; \mention{the lucky few}; relatives \\
Complemen\-tation & Selected PPs and clauses & Not characteristic of primary naming uses & Normally absent & Partitive \mention{of}-PPs; comparative \mention{than}-phrases \\
Peripheral modifiers & \mention{only the book} & \mention{even Kim} & \mention{only you}; \mention{all we who signed up} & \mention{almost every}; \mention{hardly any}; \mention{almost ten} \\
Modifier / adjunct functions & \mention{dog houses}; \mention{that day} & \mention{Canada Day}; \mention{Sunday} & Emphatic reflexives in both functions & \mention{the few people}; degree adjuncts \\
Meaning & Descriptive properties and relations & Naming in primary uses & Person, deixis, anaphora, and interrogation & Quantification, definiteness, deixis, and interrogation \\
Pro-form gender & Descriptions shape referent construal & Names identify referents construed by gender & Gender-sensitive forms & Gender-sensitive forms and constructions \\
Inflection & Number and genitive & Genitive; restricted number & Case, including genitive; reflexive forms & Number; grade; genitive; \mention{no}/\mention{none} \\
Interrogative / relative constructions & Through another constituent, e.g.\ \mention{which book} & Through another constituent & \mention{who}, \mention{whose} & \mention{which}, \mention{what}, and \mention{-ever} forms \\
Inventory & Open & Open to new names & Closed & Closed \\
\bottomrule
\end{tabular}
\end{table}

The nominal and adjectival comparisons each involve connected properties. On the noun side, quantity combines with selected complements, number transparency, and restricted dependents. On the adjective side, grade combines with degree modifiers and comparative complements. The controls in §\ref{sec:quant-controls} show why independent use alone can't decide between them. The four gradable quantifiers' limited number isn't itself a reason to discount their adjectival pattern.

The nominal connections have greater reach across the determinative system. Partitives link \mention{many} with \mention{some}; relative postmodification links \mention{few} with demonstratives and compounds. Determination and reference connect determinatives with genitive NPs and pronouns, while demonstrative number supplies a morphological parallel. These links recur across constructions whose conditions were specified separately. The gradable quantifiers participate in this wider range alongside their adjectival pattern.

I favour Noun membership because it brings the words with these recurring nominal properties into one category. Grade and degree selection still characterize a narrower subgroup. Separate D can describe the same distributions, but places a primary boundary across those connections. Determinative retains its specialization in determination and its distinctive dependent patterns within Noun, just as pronouns retain their differences from common nouns.

The broader Noun grouping still needs an account of restricted members. Their inclusion depends on how they participate in the determinative system.

\section{Restricted members and the scope of the grouping}\label{sec:articles}

The category-level argument draws on constructions that many determinatives enter. Articles lack ordinary independent argument uses, so their inclusion needs a further argument: they belong to the determinative system whose broader profile supports Noun membership. The \mention{no}/\mention{none} paradigm raises a related question about restrictions on forms.

\subsection{Articles within the determinative system}\label{sec:article-membership}

Why categorize the articles \mention{the} and \mention{a} as nouns if they can't form independent argument NPs? \mention{Every} is restricted too, while \mention{no} has the distinct independent form \mention{none} \citep[371--372, 410--411]{huddleston2002}.

\textit{CGEL}'s treatment of \mention{my} supplies a precedent for a restricted noun form \citep[470--471]{huddleston2002}. Dependent \mention{my} and independent \mention{mine} belong to one pronoun paradigm. \mention{My} remains a noun even though, in the ordinary NP construction, it requires a following nominal. The articles lack an independent argument form; their positive evidence is integration into determinative, and their nounhood depends on the category-level argument.

The articles participate in the determinative system of definiteness, quantity, and count restrictions. In \mention{the/a/this/every book}, they occupy the same Det position and contribute to the interpretation of the NP. \mention{The} permits singular, plural, and non-count targets; \mention{a} selects a singular count target and contributes individuation. Their connections extend beyond shared position, although neither has the full determinative profile \citep[368--373]{huddleston2002}.

\mention{The} also functions as a degree modifier. In \mention{the bigger the better}, it modifies comparative AdjPs, as other determinatives do in \mention{much bigger} or \mention{no better} \citep[1131--1132, 1135--1136]{huddleston2002}. This supplies a connection outside ordinary Det function while preserving its exclusion from independent argument use.

\mention{A} also enters the complex determinatives \mention{a few}, \mention{a little}, and \mention{many a}. In the last, \mention{many} contributes a large number and \mention{a} an individuating, distributive effect. \textit{CGEL} treats \mention{many a} as syntactically fixed and restricted to Det function: it doesn't establish that ordinary \mention{a} freely accepts modifiers \citep[392--394]{huddleston2002}. It does connect the article with quantificational constructions beyond \mention{a book}.

\textcite[153--158]{spinillo2004reconceptualising} instead retains \mention{the}, \mention{a}, and \mention{every} as an expanded article category and redistributes other determinatives among adjectives and pronouns \citep[194--195]{spinillo2004reconceptualising}. Her grounds include dependence on a following nominal, lack of predicative and partitive uses, and limited descriptive content. She also recognizes differences within the article trio.

\mention{Every} connects with independent \mention{each} through universal quantification and singular count selection. It also permits \mention{almost}/\mention{nearly} and occurs after genitives in \mention{her every move}. \textcite[156--158]{spinillo2004reconceptualising} recognizes these differences from \mention{the} and \mention{a}, as well as the articles' greater phonological dependence.

The exclusions of bare, predicative, and partitive uses all limit occurrence without a following nominal; they aren't three independent reasons for a primary category. Little descriptive content extends well beyond the trio. Meanwhile, \mention{a}'s quantificational constructions and \mention{every}'s modifier permissions cross the proposed article boundary. I retain the restricted members within determinative because these connections preserve a broader system without granting unrestricted use to any member.

\textcite[140--144]{spinillo2004reconceptualising} draws an analogy with verbs and prepositions whose complements can be omitted. For demonstratives and quantifiers such as \mention{some} and \mention{all}, occurrence before a nominal (\mention{some teachers}) alternates with occurrence without one (\mention{some left}). The analogy supports category continuity, but leaves the category to be established: \textit{CGEL}, Hudson, and the D-noun analysis can all preserve it.

The argument for retaining the articles is synchronic. Grammaticalization supplies background: \textcite[331--336]{lyons1999} discusses the development of definite articles from demonstratives and of articles expressing cardinality from numerals. These histories can explain changes in content and use, but don't establish the synchronic categorization.

\subsection{Dependent and independent forms}\label{sec:form-selection}

\mention{No}/\mention{none} illustrates a restriction on forms within a paradigm. \textit{CGEL} treats them as inflectional forms of one determinative: dependent \mention{no students} contrasts with independent \mention{none}, while both permit \mention{almost} \citep[389--390]{huddleston2002}. The partitive \mention{none of the students} requires the independent form; \ungram{\mention{no of the students}} is excluded. Form selection remains necessary, just as with \mention{my}/\mention{mine}.

The broader grouping preserves two distinctions: between members' permitted constructions, and between the forms selected within a paradigm.

\section{Ordinary Head, fusion, and modifier attachment}\label{sec:evidence}

The D-noun package pairs the proposed categorization with ordinary Head. This section compares that treatment with fusion in independent \mention{some}, \mention{some of the wine}, and \mention{the lucky few}. Ordinary headedness gives the determinative Head alone; fusion combines Head with a dependent function.

\subsection{The basic Head relations}\label{sec:head-relations}

A \term{nominal} (Nom) contains a head and its internal dependents, excluding an external determiner. In \mention{some apples}, the word \mention{apples} heads Nom, and that Nom heads NP. The same arrangement permits an internal modifier, as in \mention{some red apples}, without making the determining phrase part of the Nom.

Head is the function of an immediate constituent within a phrase. The \term{ultimate lexical head} is the word reached by following Head relations through the intervening phrase layers. Thus \mention{apples} is the ultimate lexical head of \mention{some apples}, although the NP's immediate Head is Nom. Both accounts also have \mention{some} as the ultimate lexical head of independent \mention{some}.

Under D-noun, a determinative can head this Nom--NP structure too. Figure~\ref{fig:some} compares the accounts for \mention{take some apples} and \mention{take some}. The dependent expression still has \mention{apples} as its ultimate head. Independent \mention{some} has ordinary Head under D-noun; in \textit{CGEL}, its DP jointly fills Det of NP and Head of Nom.

\begin{figure}[H]
\centering
\begin{minipage}[t]{.49\linewidth}\centering
\textit{CGEL}\par\smallskip
\begin{minipage}[t]{.56\linewidth}\centering
Dependent\par\smallskip
\begin{forest} nominal tree
[VP
 [{\synnode{Head}{V}}, head edge [\mention{take}]]
 [{\synnode{Obj}{NP}}
  [{\synnode{Det}{DP}}
   [{\synnode{Head}{D}}, head edge [\mention{some}]]]
  [{\synnode{Head}{Nom}}, head edge
   [{\synnode{Head}{N\textsubscript{common}}}, head edge [\mention{apples}]]]]]
\end{forest}
\end{minipage}\hfill
\begin{minipage}[t]{.42\linewidth}\centering
Independent\par\smallskip
\begin{forest} nominal tree
[VP
 [{\synnode{Head}{V}}, head edge [\mention{take}]]
 [{\synnode{Obj}{NP}}
  [{\synnode{Head}{Nom}}, head edge, before drawing tree={x+=1.5em}
   [{\synnode{Det--Head}{DP}}, no edge
    [{\synnode{Head}{D}}, head edge [\mention{some}]]]
   {\draw[-] (!uu.south) -- (); \draw[-,line width=1pt] (!u.south) -- (); }]]]
\end{forest}
\end{minipage}
\end{minipage}\hfill
\begin{minipage}[t]{.49\linewidth}\centering
\textit{D-noun analysis}\par\smallskip
\begin{minipage}[t]{.56\linewidth}\centering
Dependent\par\smallskip
\begin{forest} nominal tree
[VP
 [{\synnode{Head}{V}}, head edge [\mention{take}]]
 [{\synnode{Obj}{NP}}
  [{\synnode{Det}{NP}}
   [{\synnode{Head}{Nom}}, head edge
    [{\synnode{Head}{N\textsubscript{D}}}, head edge [\mention{some}]]]]
  [{\synnode{Head}{Nom}}, head edge
   [{\synnode{Head}{N\textsubscript{common}}}, head edge [\mention{apples}]]]]]
\end{forest}
\end{minipage}\hfill
\begin{minipage}[t]{.42\linewidth}\centering
Independent\par\smallskip
\begin{forest} nominal tree
[VP
 [{\synnode{Head}{V}}, head edge [\mention{take}]]
 [{\synnode{Obj}{NP}}
  [{\synnode{Head}{Nom}}, head edge
   [{\synnode{Head}{N\textsubscript{D}}}, head edge [\mention{some}]]]]]
\end{forest}
\end{minipage}
\end{minipage}
\caption{\mention{Take some apples} and \mention{take some} under \textit{CGEL} (first pair) and D-noun (second pair). Within each pair, the dependent use precedes the independent use. In the second diagram, DP fills Det of NP and Head of Nom. In the D-noun pair, \mention{some} has the same nominal projection in Det and Obj uses. Function labels appear above category labels. Det marks determiner, Obj object, and Det--Head fusion; the two links show the combined functions. N, D, and V mark noun, determinative, and verb; NP, DP, and VP mark their phrases. Subscripts identify noun subcategories; NP\textsubscript{D} marks a determinative-headed NP.}\label{fig:some}
\end{figure}

In this grammar, \term{inheritance} applies constraints stated for Noun to its subcategories, subject to their lexical and constructional restrictions. Ordinary headedness supplies the direct Nom--NP projection used here. A separate-D ordinary-Head grammar can license the same structure by admitting both Noun and D.

The genitive NP \mention{Kim's} already fills Det in \mention{Kim's preferences}, as Figure~\ref{fig:genitive} shows. Ordinary determinative headedness gives \mention{some} that phrase type in Det and object uses alike. Plain determinative-headed NPs, such as \mention{almost ten} in \mention{almost ten apples}, then join genitives such as \mention{Kim's} and \mention{my}. The fragment compares the remaining selectional conditions (§\ref{sec:det-uniform}).

\begin{figure}[H]
\centering
\begin{forest} nominal tree
[NP
 [{\synnode{Det}{NP[gen]}}
  [{\synnode{Head}{Nom}}, head edge
   [{\synnode{Head}{N\textsubscript{proper}}}, head edge [\mention{Kim's}]]]]
 [{\synnode{Head}{Nom}}, head edge
  [{\synnode{Head}{N\textsubscript{common}}}, head edge [\mention{preferences}]]]]
\end{forest}
\caption{The genitive NP \mention{Kim's} functions as determiner in the larger NP \mention{Kim's preferences}, whose ultimate lexical head is \mention{preferences}. The representation abstracts from the internal realization of genitive marking.}\label{fig:genitive}
\end{figure}

\subsection{Syntactic completeness and contextual interpretation}\label{sec:saturation}

With a group of people under discussion, \mention{Some left}, \mention{Many came}, and \mention{All agree} illustrate \term{syntactic completeness}: an argument expression can be complete without another overt head or determiner. Bare \mention{few} shows the same with no group under discussion: \mention{few would dare underestimate Aaron Finch's men} and \mention{few have had the opportunity to digest its contents}.\footnote{NOW: \textit{Daily Times} (Pakistan), 20 October 2022; \textit{New York Times}, 4 March 2022. Also \enquote{Everyone can have the ability to do this, but few have the durability to do it} (\textit{Los Angeles Times}, 27 September 2021). In 89 of 94 NOW lines for \mention{few would}, nothing in the visible context supplies the set. Screened lines under \texttt{corpus/independent\_argument/}.} In \mention{I'll take some}, the relevant substance or set may still be supplied by discourse or the situation. Ordinary pronouns also depend on context.

In the following web-review sentence, \mention{two different Honda models} supplies the domain for the independent object \mention{both}:\footnote{English Web Treebank, sentence \nolinkurl{reviews-083459-0002}, verified in the \href{https://github.com/UniversalDependencies/UD_English-EWT/blob/master/en_ewt-ud-train.conllu}{UD English EWT training data}. CGELBank supplies the syntactic annotation \citep{reynolds2023unified}. The original review webpage wasn't identified.}

\ea\label{ex:attested-both}
\mention{Went there yesterday: we are trying to decide between two different Honda models, so we wanted to test-drive both back to back.}
\z

Generalizing expressions such as \mention{Many are called, few are chosen} and \mention{Enough is enough} need no previously uttered common-noun phrase. Interpretation can instead depend on the situation or a generic restriction. Generic pronoun \mention{one}, as in \mention{One shouldn't judge}, provides a parallel without a required overt antecedent.

The comparison with ordinary pronouns and proper nouns is structural as well as interpretive. In \mention{She left} and \mention{Kim left}, \textit{CGEL} permits an NP headed by a nominal with no determiner. Applying that structure to \mention{Some left} preserves the same division between a complete NP and its context-dependent reference. The dependent use of \mention{some} doesn't by itself require a Det function inside every independent occurrence.

Fusion offers another representation of that complete expression. In \textcite{Payne2007}, one constituent jointly realizes Head and a dependent function, subject to structural and adjacency conditions. Contextual recovery alone doesn't decide between this arrangement and ordinary Head. Both analyses retain an overt determinative and can express the restriction supplied by context.

\subsection{Partitives and externally determined expressions}\label{sec:modification}

In \mention{some of the wine}, \mention{some} specifies a quantity drawn from an identifiable amount of wine. The NP \mention{the wine} denotes the partitive domain and is complement of \mention{of}. A pronoun-headed NP or independent genitive NP can express the domain too: \mention{many of them}, referring to previously mentioned people, or \mention{some of Kim's}, referring to Kim's apples.

Compare where the common noun occurs in \mention{some apples} and \mention{some of the wine}. In the partitive, \mention{wine} is embedded inside the \mention{of}-phrase, so it can't be the ultimate lexical head of the whole NP. The Head relations in both analyses lead instead to \mention{some}.

Figure~\ref{fig:partitive} gives the D-noun analysis, with the \mention{of}-phrase functioning as complement within Nom, as in \textit{CGEL}'s partitive tree \citep[411--412]{huddleston2002}. In \textit{CGEL}, the DP headed by \mention{some} fills Det--Head. Both analyses let the overt quantifier and domain supply a compositional interpretation. They agree on the ultimate lexical head and differ in the intervening Head relations.

\begin{figure}[H]
\centering
\begin{forest} nominal tree
[NP
 [{\synnode{Head}{Nom}}, head edge
  [{\synnode{Head}{N\textsubscript{D}}}, head edge [\mention{some}]]
  [{\synnode{Comp}{PP}}
   [{\synnode{Head}{P}}, head edge [\mention{of}]]
   [{\synnode{Comp}{NP}} [\mention{the wine}, roof]]]]]
\end{forest}
\caption{\mention{Some of the wine} under the D-noun analysis. The Head relations lead from the outer NP to \mention{some}. The common noun \mention{wine} is inside the complement PP and doesn't head the whole expression. The inner NP is abbreviated.}\label{fig:partitive}
\end{figure}

Independent \mention{few} and certain other indefinite determinatives permit definite determination. Under D-noun, \mention{few} heads an ordinary NP both alone and in \mention{the few}, where \mention{the} fills Det. Adding an optional adjective gives \mention{the lucky few}, with the same Head and dependent functions as \mention{the lucky survivors}.\footnote{\mention{The lucky few} is well attested as an independent argument with a relative: \enquote{Are you one of the lucky few who found powdered infant formula?} (\textit{National Observer}, 3 October 2022); \enquote{is among the lucky few who already have that permanent freedom} (\textit{Business Insider}, 11 February 2021); \enquote{Obialor is one of the lucky few who was able to reclaim the moment} (YNaija, 11 September 2020); likewise \enquote{I am of the privileged few who can afford to take time off} (\textit{The Guardian}, 1 September 2022). All from NOW, 17 September 2026; screened lines under \texttt{corpus/independent\_argument/}.}

\textit{CGEL} explicitly permits determinatives used as internal modifiers to fuse with Head, as in \mention{the other two} and \mention{these few here} \citep[415--416]{huddleston2002}. Its analysis of \mention{the few mistakes} assigns \mention{the} to Det and \mention{few} to Mod \citep[392]{huddleston2002}. That dependent use supplies the counterpart for a Mod--Head analysis of independent \mention{few} after an external determiner.

Figure~\ref{fig:few} compares the resulting analyses of \mention{the lucky few}. Both have an overt Det and a Nom modified by \mention{lucky}. In the D-noun analysis, the lexical word \mention{few} directly heads Nom. In the fusion analysis, its DP fills Mod--Head. The competing structure thus uses machinery already needed for adjectival fusion; it requires neither a silent noun nor conversion.

\begin{figure}[H]
\centering
\begin{minipage}{.48\linewidth}\centering
\begin{forest} nominal tree
[NP
 [{\synnode{Det}{NP\textsubscript{D}}} [\mention{the}, roof]]
 [{\synnode{Head}{Nom}}, head edge
  [{\synnode{Mod}{AdjP}} [\mention{lucky}]]
  [{\synnode{Head}{N\textsubscript{D}}}, head edge [\mention{few}]]]]
\end{forest}
\end{minipage}\hfill
\begin{minipage}{.48\linewidth}\centering
\begin{forest} nominal tree
[NP
 [{\synnode{Det}{DP}} [\mention{the}, roof]]
 [{\synnode{Head}{Nom}}, head edge, before drawing tree={x/.option=!21.x}
  [{\synnode{Mod}{AdjP}} [\mention{lucky}]]
  [{\synnode{Head}{Nom}}, head edge, before drawing tree={x+=1.5em}
   [{\synnode{Mod--Head}{DP}}, no edge
    [{\synnode{Head}{D}}, head edge [\mention{few}]]]
   {\draw[-] (!uu.south) -- (); \draw[-,line width=1pt] (!u.south) -- (); }]]]
\end{forest}
\end{minipage}
\caption{\mention{The lucky few} with ordinary Head (left) and Mod--Head fusion (right). The determiner is separately realized in both. Internal structure within the article phrase and AdjP is suppressed.}\label{fig:few}
\end{figure}

Table~\ref{tab:matched-heads} compares the Head relations across five constructions; §\ref{sec:det-uniform} supplies the constraints under both taxonomies.

\begin{table}[H]
\centering\small
\caption{Head and dependent functions under matched analyses. Each independent expression is an NP; the phrase headed by \mention{the} fills Det wherever it appears. Form labels abbreviate the corresponding constituents shown in the trees. The ordinary-Head column applies to both taxonomies permitting ordinary determinative heads.}\label{tab:matched-heads}
\begin{tabular}{>{\raggedright\arraybackslash}p{2.8cm}>{\raggedright\arraybackslash}p{4.5cm}>{\raggedright\arraybackslash}p{4.5cm}}
\toprule
Expression & Separate D with fusion & Ordinary determinative Head \\
\midrule
\mention{few survivors} & \mention{few}: Det; \mention{survivors}: Head & Same functions \\
Independent \mention{few} & \mention{few}: Det--Head & \mention{few}: Head \\
\mention{the few} & \mention{few}: Mod--Head & \mention{few}: Head \\
\mention{the lucky few} & \mention{few}: Mod--Head; \mention{lucky}: Mod & \mention{few}: Head; \mention{lucky}: Mod \\
\mention{the idle rich} & \mention{rich}: Mod--Head; \mention{idle}: Mod & Same fusion and modifier functions \\
\bottomrule
\end{tabular}
\end{table}

Fusion groups \mention{the lucky few} with \mention{the idle rich}; ordinary headedness groups it with \mention{the lucky survivors} and unifies bare and externally determined \mention{few}. The choice is between sharing fusion across D and adjective, and sharing ordinary Head structure across determinative constructions.

Reducing fusion's applications can simplify this description even when fusion remains available for adjectives. The advantage is the uniform treatment of \mention{few}; its value depends on any additional conditions or lost generalizations elsewhere.\footnote{Compare the theory of second best in \textcite[11--12]{lipseylancaster1956secondbest}: under a constraint preventing an optimum, satisfying more optimality conditions needn't improve the outcome. The analogy concerns interactions among grammatical choices, not a formal optimum for the grammar.} Section~\ref{sec:costs} compares those costs while holding the lexical restrictions fixed.

The constructed \mention{the remaining three} extends the modifier pattern to cardinals. It doesn't independently decide whether \mention{three} is determinative or has a common-noun use: \textcite{reynolds2026numerals} argues for both uses of cardinals. Section~\ref{sec:costs} treats their unification as a consequence conditional on that analysis.

\subsection{Modifier selection and attachment}\label{sec:payne}

The selection and ordering contrast in §\ref{sec:adjectival-profile} guides the attachment analysis. Approximative and focusing modifiers have established NP parallels; the degree series selects the four gradable quantifiers. Both ordinary-Head accounts can represent that difference by distinguishing NP-peripheral modification from modification inside Nom.

Peripheral modifiers attach outside an NP's internal dependents. Compare \mention{only you}, \mention{for almost my entire life}, and \mention{Usually a careful driver, Anne found her mind wandering}. The modifiers precede any determiner belonging to the NP they modify. \textit{CGEL} already assigns them NP-level attachment \citep[430--431]{huddleston2002}. Focusing \mention{even} and \mention{just}, and approximatives such as \mention{nearly}, \mention{hardly}, \mention{virtually}, and \mention{practically}, likewise have NP uses, subject to their own scope and selectional restrictions.

In \mention{almost every experienced teacher}, \mention{almost} belongs with \mention{every}, while \mention{experienced} modifies \mention{teacher}. Figure~\ref{fig:every} places \mention{almost} at the periphery of the NP headed by \mention{every}. That smaller NP functions as Det in the larger NP, whose ultimate head is \mention{teacher}.

\begin{figure}[H]
\centering
\begin{forest} nominal tree
[NP
 [{\synnode{Det}{NP\textsubscript{D}}}
  [{\synnode{Mod}{AdvP}} [\mention{almost}]]
  [{\synnode{Head}{NP\textsubscript{D}}}, head edge
   [{\synnode{Head}{Nom}}, head edge
    [{\synnode{Head}{N\textsubscript{D}}}, head edge [\mention{every}]]]]]
 [{\synnode{Head}{Nom}}, head edge
  [{\synnode{Mod}{AdjP}} [\mention{experienced}]]
  [{\synnode{Head}{N\textsubscript{common}}}, head edge [\mention{teacher}]]]]
\end{forest}
\caption{Peripheral \mention{almost} modifies the smaller NP headed by \mention{every}; the resulting \mention{almost every} NP functions as Det. The adjective \mention{experienced} modifies the common-noun nominal. Internal structure within the one-word modifier phrases is suppressed.}\label{fig:every}
\end{figure}

\textcite[40--42]{payne2010} defend keeping \mention{few}, \mention{any}, and related forms in one lexical category across dependent and independent uses. In \mention{hardly any money} and independent \mention{hardly any}, \mention{hardly} remains an adverb and \mention{any} retains its category. The D-noun analysis preserves that continuity. It also preserves their contrast between adverbial premodification, as in \mention{almost anybody}, and adjectival postmodification, as in \mention{nothing absolute} \citep[41--42]{payne2010}.

The proposed analysis gives adverbs such as \mention{very} and \mention{too} internal attachment in dependent and independent uses alike. In \mention{the very few people}, \mention{very} modifies the Nom headed by \mention{few}; the resulting \mention{very few} NP functions as Mod in the larger NP. In \mention{the very few who objected}, \mention{few} remains ordinary Head and \mention{very} remains internal. The degree relation is unchanged, and one internal permission covers both constructions. The same analysis applies to \mention{the very many}.

NP degree modifiers provide a further connection. Determinatives permit determinative premodifiers such as \mention{this} in \mention{this much} \citep[393]{huddleston2002}. Comparative determinatives also take established NP modifiers: \mention{a lot} in \mention{a lot fewer}, alongside \textit{CGEL}'s \mention{a lot more than fifty} \citep[432]{huddleston2002}. These modify the quantity expressed by the head. Under the proposed analysis, they're internal NP modifiers within its nominal projection, with permissions specific to the head and construction.

\textcite[42--47]{payne2010} establish that adverbs can postmodify common nouns, as in the constructed \mention{the changes globally to the climate}. They also analyse \mention{almost} as modifying the attributive nominal \mention{textbook} in \mention{an almost textbook case} \citep[p.~75, n.~3]{payne2010}. Adverbial modification therefore isn't categorically excluded from Noun, even internally. Its broader availability with determinatives remains a difference of distribution within the proposed category.

Approximatives also need internal attachment where they follow an external determiner on an independent cardinal: \mention{the almost thirty who came}. The internal AdvP permissions at issue are thus the degree series on \mention{few}, \mention{many}, \mention{much}, and \mention{little} across dependent and independent uses, and approximatives on externally determined independent cardinals.

\subsection{Compounds and their modifier domains}\label{sec:compounds}

A more demanding modifier comparison is \mention{hardly anyone present}. \textcite[581--583]{Payne2007} assign \mention{hardly} to DP structure and \mention{present} to nominal structure. The compound \mention{anyone} takes the premodifiers of its determinative base \mention{any}: compare \mention{hardly any writer present}. The adjective realizes a specialized \term{restrictor} function, restricted to post-head position and non-recursive. Figure~\ref{fig:compound} contrasts this account with an ordinary-Head analysis.

\begin{figure}[H]
\centering
\begin{minipage}{.48\linewidth}\centering
\begin{forest} nominal tree
[NP
 [{\synnode{Mod}{AdvP}} [\mention{hardly}]]
 [{\synnode{Head}{NP}}, head edge
  [{\synnode{Head}{Nom}}, head edge
   [{\synnode{Head}{N\textsubscript{D}}}, head edge [\mention{anyone}]]
   [{\synnode{Mod}{AdjP}} [\mention{present}]]]]]
\end{forest}
\end{minipage}\hfill
\begin{minipage}{.48\linewidth}\centering
\begin{forest} nominal tree
[NP, before drawing tree={x/.option=!11.x}
 [{\synnode{Head}{Nom}}, head edge, before drawing tree={x+=4em}
  [{\synnode{Det--Head}{DP}}, no edge
   [{\synnode{Mod}{AdvP}} [\mention{hardly}]]
   [{\synnode{Head}{D}}, head edge [\mention{anyone}]]]
  {\draw[-] (!uu.south) -- (); \draw[-,line width=1pt] (!u.south) -- (); }
  [{\synnode{Mod}{AdjP}}, before drawing tree={x/.option=!u.x} [\mention{present}, before drawing tree={x/.option=!u.x}]]]]
\end{forest}
\end{minipage}
\caption{\mention{Hardly anyone present} with ordinary Head (left) and fusion (right). On the left, \mention{hardly} is peripheral to NP and \mention{present} is internal to Nom. Both post-head Mod positions realize the specialized restrictor. In the fusion tree, DP fills Det of NP and Head of Nom, following \textcite[p.~582, (13d)]{Payne2007}. Modifier-phrase interiors are suppressed.}\label{fig:compound}
\end{figure}

In the ordinary-Head tree, \mention{anyone} inherits nominal projection from Noun and its premodifier permissions from its determinative base. Peripheral \mention{hardly} follows the approximative pattern and its NP parallels in §\ref{sec:payne}. The compound construction independently excludes external determination and supplies the internal post-head restrictor permission, with the same position and non-recursion conditions as the fusion account. It doesn't license a corresponding pre-head adjective. Both exclusions hold of the compound as a determinative. The compounds also have a secondary set-denoting use, in which they take an article, a pre-head adjective, and plural inflection: lexicalized \mention{a special someone} and \mention{a certain someone}, and productively \mention{an anonymous someone}, \mention{the responsible someone}, \mention{a furry something}, and \mention{such an anyone}.\footnote{COCA: \mention{special someone} 282 tokens and \mention{certain someone} 120 (List counts, 17 September 2026); \enquote{if an anonymous someone with very little solid evidence on his/her claims can get basically every major media outlet} (\textit{Mashable}, 2017); \enquote{The responsible someone has to collect the rents} (\textit{Confrontation}, 2014); \enquote{a furry something swinging indelicately from its mouth. The something plopped wetly into her bowl} (\textit{Analog}, 2002), where the next sentence takes \mention{the}; \enquote{Not anyone. And if I cld think of such an anyone, I would go out there} (\textit{Bones}, 2010, transcript spelling). Screened lines under \texttt{corpus/exclusion\_tests/}.} \textit{CGEL} records \mention{a nobody} and \mention{a little something} as nouns formed by conversion \citep[423, n.~43]{huddleston2002}. Under the D-noun analysis no category change is involved: the use is the one \textit{CGEL} gives proper nouns in \mention{an Ophelia} and \mention{the Smiths}, where the name comes to \enquote{denote the set of people bearing this name} and \enquote{as a set-denoting noun} takes determiners and restrictive modifiers \citep[521]{huddleston2002}. In both subcategories the article and the pre-head adjective follow from what the noun now denotes, a kind or an individual of the quantifier's description, and quantificational force is gone: \mention{an anybody} is a person of no particular standing. The quantifier and its nominal echo can share a text: \enquote{anybody could have competed. I qualified as an anybody} (\textit{Skiing}, 1992, COCA). One process thus covers two subcategories where \textit{CGEL} needs conversion for one and a secondary use for the other. Locative compounds before a noun (\mention{an anywhere operation}) are the ordinary noun-as-modifier use and take their article from the head noun.\footnote{Whether the set-denoting use still admits the post-head restrictor (\mention{a certain someone special}) is a judgment I haven't tested; the corpus lines show article and pre-head adjective without it.}

The relative clauses in §\ref{sec:independent} are ordinary postmodifiers, distinct from the specialized restrictor. When both occur, the relative follows it: \mention{something useful that I found}. \textit{CGEL} explicitly gives compounds the common noun's range of ordinary postmodifiers while preserving this ordering condition \citep[423]{huddleston2002}.

The same division applies to \mention{someone} and \mention{everybody}: the compound construction licenses the post-head restrictor, while the determinative base supplies its premodifier permissions. The example \mention{something reliable and good looking} illustrates the post-head pattern (§\ref{sec:independent}).

Both accounts separate the modifier domains structurally: fusion uses the DP--Nom boundary, and ordinary headedness uses the NP--Nom boundary. Their constituent groupings differ: ordinary Head groups \mention{anyone present} in the inner NP, whereas fusion groups \mention{hardly anyone} in DP.

\subsection{Independent genitives}\label{sec:genitive-head}

\textit{CGEL} also uses fusion with nouns. Anaphoric \mention{mine}, understood as \enquote*{my car}, combines the possessive relation with an understood nominal description and receives a fused analysis. Predicative \mention{mine}, as in \mention{it's mine}, can instead mean \enquote*{belongs to me}: it expresses the relation directly, without an understood nominal head \citep[410--411, 470--471]{huddleston2002}. A suitable context may permit either reading.

The proposed analysis gives \mention{mine} ordinary Head in both anaphoric and predicative uses. Its genitive form expresses a relation to the speaker; context can supply the understood nominal description, as it can with independent \mention{few}. The \mention{my}/\mention{mine} alternation remains a form-selection condition, parallel to \mention{no}/\mention{none}. This extends ordinary headedness to an existing noun subcategory.

The genitive's person identifies the possessor; the independent NP's agreement reflects the possessed entity or entities. Compare constructed \mention{Mine is ready} (\enquote*{my contribution}) and \mention{Mine are ready} (\enquote*{my slides}). Ordinary Head therefore needs a constructional agreement condition. This extension isn't counted as a gain from determinative membership within Noun.

\section{The matched fragment and comparative costs}\label{sec:economy}

A grammatical \term{fragment} states conditions on a specified range of constructions. Each account assigns categories and Head relations to the same expressions. The conditions specify permitted functions, dependents, and combinations of forms. They describe structures directly, following the constraint-based approach outlined by \textcite{pullum2020theorizing}. Comparing their content and reuse makes the accounts' commitments explicit; counting written statements wouldn't measure grammatical economy.

\subsection{Scope and use permissions}\label{sec:fragment}

The fragment covers determination, internal nominal modification, ordinary subject, object, and complement-of-preposition uses, partitives, and the compound comparisons in §\ref{sec:evidence}. It admits relative postmodifiers without analysing their internal grammar. Predication, interrogative clause structure, and modification outside NP structure lie beyond its scope. Degree uses enter the wider comparison in §\ref{sec:enough}.

Table~\ref{tab:permissions} separates uses before another nominal from independent argument uses. These are lexical \term{use permissions}, with separate conditions selecting the appropriate form. A permission applies to an occurrence in the specified construction. Intermediate Head relations within its projection don't require a further argument permission.

\begin{table}[H]
\centering\small
\caption{Selected permissions and form restrictions in the fragment. Det and Mod concern uses before another nominal; Mod is internal. A dash marks absence of the stated permission. Unstressed \mention{some} has the target restriction shown.}\label{tab:permissions}
\setlength{\tabcolsep}{4pt}
\begin{tabular}{lcccc>{\raggedright\arraybackslash}p{4.7cm}}
\toprule
Form & Det & Mod & Subj & Obj / CompP & Target before another nominal \\
\midrule
\mention{the} & yes & -- & -- & -- & No count or number restriction \\
\mention{a} & yes & -- & -- & -- & Singular count \\
\mention{every} & yes & yes & -- & -- & Singular count; Mod after genitive Det \\
\mention{some} & yes & -- & yes & yes & Plural count or non-count \\
\mention{few} & yes & yes & yes & yes & Plural count \\
\mention{no} & yes & -- & -- & -- & No count or number restriction \\
\mention{none} & -- & -- & yes & yes & Not applicable \\
\mention{my} & yes & -- & -- & -- & No count or number restriction \\
\mention{mine} & -- & -- & yes & yes & Not applicable \\
\mention{she} & -- & -- & yes & -- & Not applicable \\
\mention{her} (plain) & -- & -- & -- & yes & Not applicable \\
\bottomrule
\end{tabular}
\end{table}

Subj and Obj denote subject and object; CompP denotes complement of a preposition. The \mention{no}/\mention{none} and \mention{my}/\mention{mine} pairs distinguish forms within a paradigm (§\ref{sec:form-selection}). Plain \mention{she}/\mention{her} distinguishes subject from object and prepositional-complement uses. Quotation, metalinguistic naming, and subordinate-clause uses of dependent genitives fall outside the fragment.

In \mention{some apples}, \mention{some}'s phrase satisfies its Det permission and plural-count target selection. In \mention{Some left}, it satisfies its subject permission. \mention{Every apple} passes the dependent checks; independent \ungram{\mention{Every arrived}} fails its use permission.

Target restrictions apply to the nominal being determined or modified, including a nominal headed by an independent determinative. Thus \mention{the} permits plural \mention{few} in \mention{the few}. Determinative-headed phrases also bear number and count properties: \mention{few} is plural count, while \mention{some} takes its interpretation from the construction and domain. The modifier \mention{this} in \mention{this much} expresses degree; its permission doesn't follow from demonstrative Det selection.

\textcite[7--8]{hudson2004determiners} distinguishes dependency from external headedness. He argues that determiner and common noun depend on each other, although only one connects the phrase to its surroundings. His temporal adjunct evidence supports common-noun headedness: \mention{I saw him that day} is possible, whereas \ungram{\mention{I saw him that point in time}} isn't, despite the similar temporal meanings \citep[10--12]{hudson2004determiners}. The lexical noun matters, as well as the determiner restrictions.

Reciprocal selection is compatible with this headedness. A singular count noun normally requires determination, while \mention{every} requires another nominal and \mention{some} doesn't. The noun can remain Head while both constituents impose conditions.

In the external uses listed above, an expression has to pass three checks: permission for its use, selection of its form, and satisfaction of target and dependent conditions. For a phrase occurrence $x$ with external function $f$ in a surrounding structure $s$, the shared use condition collects these checks:

\[
\operatorname{admissible}(x,f,s)
\quad\Longleftrightarrow\quad
f\in U_{\ell(x)}\ \land\ \operatorname{Form}(x,f,s)\ \land\ C(x,f,s).
\]

Here $\ell(x)$ identifies the lexical head's lexeme and $U_{\ell(x)}$ its use permissions. $\operatorname{Form}$ checks the selected form, including independent \mention{none} in \mention{none of the students}. $C$ checks the occurrence's dependents, its target where relevant, and the other constructional conditions. The structure $s$ includes both $x$ and its surroundings; for \mention{some}'s phrase in \mention{some apples}, it includes the target \mention{apples}.

\subsection{Phrase categories and structural constraints}\label{sec:det-uniform}

In \mention{some apples} and \mention{Kim's apples}, the determining phrases share the NP category under D-noun with ordinary Head. Separate D can permit the same projection. For D-noun with fusion, I retain that NP category and add the combined function relations in independent uses. \textit{CGEL} instead distinguishes DP from genitive NP. It also admits restricted plain-case NPs and PPs as determiners, as in \mention{what size hat} and \mention{over thirty ties} \citep[Ch.~5, §4]{huddleston2002}.

I carry this PP-determiner provision over to the three NP-projecting implementations for the present comparison. Its inclusion in the alternatives is an extrapolation.

\[
\begin{aligned}
\text{\textit{CGEL}:} &\quad \mathrm{Det}:\{\mathrm{DP},\mathrm{NP},\mathrm{PP}\} \\
\text{NP-projecting implementations here:} &\quad \mathrm{Det}:\{\mathrm{NP},\mathrm{PP}\}
\end{aligned}
\]

NP projection consolidates the principal phrase types in Det function. Determinative-headed, genitive, and other licensed NPs still require separate identification. All four implementations retain at most one Det per NP and the same definiteness and target-selection conditions. In the separate-D ordinary-Head account, an NP's ultimate lexical head can be Noun or D.

The ordinary-Head accounts share three structural conditions. First, an NP core has a Nom as Head and at most one Det. Second, a nominal core has a lexical Head and only the dependents permitted for that head in its construction. Third, a peripheral modifier combines with an NP whose Head relation continues to an NP, as in Figure~\ref{fig:every}. This keeps peripheral modification outside the core's Det and Nom dependents.

Where a word occurrence $h$ directly fills Head in Nom, let $\operatorname{Cat}(h)$ record its primary lexical category. Subcategory distinctions remain available. The two taxonomies impose different category conditions:

\[
\begin{aligned}
\text{D-noun, ordinary Head:} &\quad \operatorname{Cat}(h)=\mathrm{N} \\
\text{Separate D, ordinary Head:} &\quad \operatorname{Cat}(h)\in\{\mathrm{N},\mathrm{D}\}.
\end{aligned}
\]

Under D-noun, determinatives satisfy the first condition as members of Noun. Separate D admits them alongside Noun. Both permit the same nominal layers, selected complements, and ordered modifiers. Their agreement on phrase structure leaves the lexical grouping open to the profile argument.

The dependent conditions distinguish uses before another nominal from independent uses. In Det or internal Mod function before a target, ordinary \mention{some} and \mention{few} exclude external determination, partitive complements, and relative postmodifiers within their own phrase. Independent \mention{some} and \mention{few} permit selected partitives and relatives; independent \mention{few} also permits definite determination and adjectival premodification, as in \mention{the lucky few}. These permissions don't transfer to \mention{every} or the articles.

Degree-modifier selection applies across the relevant dependent and independent uses. It includes AdvPs such as \mention{very}, determinative-headed modifiers such as \mention{this} in \mention{this much}, and established NPs such as \mention{a lot} in \mention{a lot fewer}. Comparative complements, as in \mention{more than ten}, are licensed separately from partitive \mention{of}-phrases. Neither the category NP nor a general modifier slot licenses arbitrary adjectives or unrestricted complements.

Postmodification permits zero or more ordered dependents, subject to the construction. The compounds permit at most one specialized post-head restrictor followed by ordinary postmodifiers: \mention{useful} precedes \mention{that I found} in \mention{something useful that I found}. They exclude external determination and a corresponding pre-head adjective as determinatives; the set-denoting use of §\ref{sec:compounds} (\mention{a special someone}) falls outside the construction. Their determinative bases supply the premodifier restrictions (§\ref{sec:compounds}).

The approximatives retain the two attachment sites motivated in §\ref{sec:payne}: peripheral NP attachment in \mention{almost every} and internal attachment in externally determined \mention{the almost thirty who came}. Both ordinary-Head accounts require this distinction.

Complex \mention{a few} and \mention{a little} have construction-specific conditions, including the position of internal degree modifiers in \mention{a very few}. Their initial \mention{a} isn't an ordinary Det selecting a plural \mention{few} target. The fragment records these combinations as complex determinatives and leaves their internal analysis beyond the position contrast unspecified. Fixed \mention{many a} retains its restriction to Det function (§\ref{sec:article-membership}).

The constraints also distinguish the uses of \mention{few}. In \mention{the few people}, its phrase is an internal Mod of the nominal headed by \mention{people}; in \mention{the lucky few}, \mention{few} heads the independently determined NP.

Each phrase's use permission follows its own lexical head. In \mention{the apple}, the smaller article phrase has Det permission; the outer NP has argument permission through \mention{apple}, whose determination requirement is satisfied. Bare \ungram{\mention{Book arrived}} fails that requirement, while \mention{Books arrived} passes. The article's exclusion from argument use remains a separate condition.

The fusion accounts retain these use, form, selection, and ordering conditions but assign different Head relations. Let $p$ be the independent NP, $m$ its Head Nom, and $x$ the shared determinative phrase. Write $\operatorname{Head}(p,m)$ for \enquote*{$m$ fills Head in $p$}, and likewise for Det and Mod. The two configurations satisfy the following relations:

\[
\begin{aligned}
\text{Det--Head:}\quad
 &\operatorname{Head}(p,m),\quad \operatorname{Det}(p,x),\quad \operatorname{Head}(m,x); \\
\text{Mod--Head:}\quad
 &\operatorname{Head}(p,m),\quad \operatorname{Mod}(m,x),\\
 &\operatorname{Head}(m,n),\quad \operatorname{Head}(n,x),\qquad \operatorname{Cat}(n)=\mathrm{Nom}.
\end{aligned}
\]

Det--Head fills the NP's Det function, excluding a second Det. Mod--Head leaves that function available for an external determiner. Its additional Nom $n$ matches Figure~\ref{fig:few}: the determinative phrase fills Mod in $m$ and Head in $n$. The ordinary adjective \mention{lucky} separately modifies $m$. These are simultaneous relations on a shared constituent, with no sequence of fusion operations assumed.

Under separate D with fusion, $x$ is a DP whose lexical Head is D. Selected premodifiers and comparative complements belong within that DP; the independent partitive complement, compound restrictor, and ordinary postmodifiers belong to nominal structure outside it, as in Figure~\ref{fig:compound}. The permissions for the complete independent NP still exclude articles and \mention{every}. Possessing a dependent Det use doesn't itself license an independent fused use.

Under D-noun with fusion, $x$ is the NP projected by the determinative noun. The same joint relations connect it to the outer nominal structure. In this construction, the inner NP excludes external determination, partitives, restrictors, and relatives: any licensed independent-use dependents occur outside it, as in the separate-D fusion account. These restrictions need stating separately from the limit of one fusion configuration around the determinative's projection.

Ordinary headedness uses the determinative's NP directly for bare, partitive, and externally determined expressions. In the fusion accounts, its DP or NP occupies the further nominal structure just specified. Both ordinary-Head accounts share structures. Under D-noun, determinatives meet the shared category condition through Noun membership; separate D admits them alongside Noun. Section~\ref{sec:overall-assessment} assesses this relation between the taxonomy and the shared structure.

Nesting determinative inside pronoun would add pronoun to the inheritance path from Noun without changing this fragment's judgments. Applying Hudson's nesting to the full determinative inventory used here is an extrapolation; §\ref{sec:inheritance} discusses the inventory differences.


\subsection{Degree uses outside noun phrases}\label{sec:enough}

When \mention{enough} modifies \mention{good} in \mention{good enough}, the three NP-projecting implementations analyse it as an NP modifier of an adjective. This would be a cost if NPs couldn't occur in that function, or if their distribution differed systematically from that of degree determinatives. The comparison therefore includes the degree uses of established NPs. Attributive uses are the crucial test, because \textit{CGEL} restricts NP modifiers there.

\mention{Enough} tests both position and external function. It precedes a nominal in \mention{enough money} and can follow one in \mention{money enough}; post-head \mention{enough} can't itself be premodified, as shown by \ungram{\mention{money almost enough}} \citep[396--397, 445]{huddleston2002}.

NP modification inside an AdjP also permits post-head position. In the constructed examples in (\ref{ex:temporal-adjp}), I analyse \mention{every Tuesday}, \mention{this week}, and \mention{next week} as temporal NP modifiers within the bracketed AdjPs. The separate degree modifiers and temporal specifications in the coordinated example make that attachment clear. These NPs express time, while \mention{enough} expresses degree; the shared position leaves their selectional differences intact.

\ea\label{ex:temporal-adjp}
\ea \mention{Applicants} [\mention{completely free every Tuesday}] \mention{will be interviewed next month.}
\ex \mention{Anyone} [\mention{very busy this week}] \mention{but} [\mention{completely free next week}] \mention{should contact me today.}
\z\z

\mention{Enough} also modifies adjectives, adverbs, verbs, and some PPs: \mention{good enough}, \mention{quickly enough}, \mention{I hadn't prepared enough}, and \mention{enough in control} \citep[396--397]{huddleston2002}. The degree determinatives \mention{much} and \mention{little}, and \mention{no}/\mention{none}, likewise have uses outside NP structure \citep[390, 395--397]{huddleston2002}.\footnote{\mention{Both}, \mention{either}, and \mention{neither} also serve as markers of coordination \citep[1305, 1308]{huddleston2002}. Their noun categorization would preserve that further category--function combination; the present fragment doesn't analyse coordination.}

NPs already modify adjectives in \mention{three years old}, \mention{a great deal smaller}, and \mention{plenty big enough}. \textit{CGEL} also explicitly treats \mention{lots better} and \mention{heaps worse} as containing quantificational NP modifiers \citep[549--550]{huddleston2002}. These uses extend the connection with the quantificational common nouns compared in §\ref{sec:quant-nouns}.

\textit{CGEL} contrasts predicative \mention{a great deal better} with \ungram{\mention{some a great deal better proposals}}, against permitted \mention{some much better proposals} \citep[551--552]{huddleston2002}. The excluded example also juxtaposes \mention{some} and \mention{a}; it doesn't isolate NP category as the cause. \textit{CGEL} permits \mention{She's a lot better player than me}, analysing the inner article as lost after the outer one \citep[p.~552, n.~8]{huddleston2002}.

Informal attributive examples with nominal \mention{heaps} and \mention{lots} further qualify a categorical exclusion.\footnote{\mention{Some heaps better photo's}: \href{https://www.aulro.com/afvb/motor-cycling/61858-my-track-day-phillip-island-motogp-track-post798794.html}{AULRO, moose, 19 August 2008, post 2}; original spelling retained. \mention{Some lots better ones}: \href{https://www.uberpeople.net/threads/let-s-play-would-you-have-taken-this-trip.395130/}{Uber Drivers Forum}, verified in an indexed excerpt; author and date remain unverified. These are informal attestations; any register restriction remains to be established.} They don't establish the same distribution as degree determinatives. The lexical, constructional, and register conditions remain a descriptive question for all four accounts. If degree determinatives systematically permit attributive uses excluded for established nominal modifiers, that difference would favour DP projection. It would bear on phrase structure, independently of determinatives' membership within Noun.

\subsection{Comparison of the four implementations}\label{sec:costs}

Table~\ref{tab:economy} compares the four implementations. It distinguishes Head configurations, modifier permissions, and taxonomic consequences. The inventory, constructions, and readings remain fixed.

\begin{table}[H]
\centering\small
\caption{Four implementations compared. All retain the lexical and constructional restrictions and adjectival fusion. Extending ordinary Head to independent genitives adds a constructional agreement condition (§\ref{sec:genitive-head}). Cardinal and \mention{one} groupings are consequences of the categorization.}\label{tab:economy}
\setlength{\tabcolsep}{3pt}
\begin{tabular}{>{\raggedright\arraybackslash}p{2.5cm}>{\raggedright\arraybackslash}p{3.6cm}>{\raggedright\arraybackslash}p{3.6cm}>{\raggedright\arraybackslash}p{2.65cm}}
\toprule
Account & Independent determinatives & Modifier permissions & Taxonomic consequences \\
\midrule
D-noun, ordinary Head & Nom--NP with Head alone across bare, partitive, and externally determined uses. & Peripheral NP premodifiers; head-specific internal modifiers. Degree determinatives project NP. & D within Noun. Cardinal uses and the three \mention{one} lexemes are grouped there. \\
Separate D, ordinary Head & Same structures, with Nom admitting either N or D as lexical Head. & Same internal and peripheral permissions and degree projection as the proposed package. & D remains separate from nominal counterparts. \\
Separate D, fused Head & DP fills Det--Head or Mod--Head in nominal structure. & DP--Nom separates premodifiers and postmodifiers. Degree determinatives project DP. & Same primary-category divisions as separate D with ordinary Head. \\
D-noun, fused Head & A projected NP fills Det--Head or Mod--Head in additional nominal structure. & Premodifiers within the fused NP; postmodifiers in the outer Nom. Degree determinatives project NP. & Same Noun grouping as the proposed package. \\
\bottomrule
\end{tabular}
\end{table}

The articles' restrictions, \mention{no}/\mention{none} form selection, and the compound restrictor conditions are already required in the separate-D grammar. Retaining them isn't an added cost of D-noun. Likewise, differences among existing noun subcategories in determination and modification remain in place (Table~\ref{tab:existing}).

Both ordinary-Head accounts distinguish peripheral NP modifiers from the internal degree and cardinal approximative permissions (§\ref{sec:payne}). For compounds, this NP--Nom boundary preserves the modifier-domain contrast that fusion expresses through its inner phrase and outer Nom (§\ref{sec:compounds}). The stated selection and attachment conditions remain necessary under both treatments.

In the ordinary-Head accounts, \mention{few} heads Nom and Nom heads NP across bare, partitive, and externally determined uses. In the fusion accounts, its DP or NP fills Det--Head in bare and partitive uses and Mod--Head with an external determiner (Table~\ref{tab:matched-heads}). In these matched examples, all have Nom as the NP's immediate Head. The contrast concerns the full Head chain and the additional fused functions.

I favour sharing the direct Head chain across determinative uses and with established nouns. The selection and attachment conditions preserve the modifier-domain contrast. Fusion instead shares a construction with adjectives. Both cover the matched cases with the same lexical restrictions; the trade-off concerns which structures the grammar shares. The wider profile comparison supplies the reason for preferring the nominal grouping.

\textcite{reynolds2026numerals} distinguishes determinative numerals such as \mention{ten} in \mention{ten men}, proper-noun cases such as \mention{10} in \mention{Room 10}, and common-noun numerals such as \mention{tens} in \mention{tens of pens}. Under the D-noun categorization, all fall within primary Noun, retaining their subcategory and constructional differences. Ordinals remain adjectives, and complex numeral phrases remain distinct from single lexemes.

\textcite[797--798]{payne2013anaphoric} distinguish three lexemes spelled \mention{one}: determinative, anaphoric common noun, and generic personal pronoun. These remain distinct under the D-noun categorization, but all belong within primary Noun.

Complex cardinals also separate category from function. In \mention{two hundred books}, the whole \mention{two hundred} fills Det; internally, \mention{two} modifies the magnitude head \mention{hundred} \citep[§4]{reynolds2026numerals}. In \mention{these two hundred books}, \mention{these} fills Det and \mention{two hundred} is an internal modifier. One Det function doesn't entail a limit of one determinative lexeme per NP.

Both D-noun implementations share this grouping of cardinals and the three \mention{one} lexemes.

\subsection{Overall assessment}\label{sec:overall-assessment}

The D-noun package combines the two preferences. The lexical comparison favours Noun membership because nominal connections recur across more of the determinative system (§\ref{sec:membership}). The structural comparison favours the direct Head chain across the licensed uses (§\ref{sec:costs}). Within the proposed grammar, determinatives then meet the same Noun condition on lexical Heads in Nom as the existing subcategories (§\ref{sec:det-uniform}). That category condition expresses the grouping supported by the profile.

Separate D with ordinary Head licenses the same structures by admitting N or D. The advantage claimed for D-noun is the correspondence between the lexical grouping and this shared grammar. Counting fewer primary categories or fewer symbols in the Head condition wouldn't establish that correspondence.

The grouping also has a consequence that counting categories doesn't. Compound determinatives and proper nouns share a secondary set-denoting use in which the word takes an article, a pre-head adjective, and plural inflection (§\ref{sec:compounds}): \mention{a special someone} and \mention{such an anyone} beside \mention{an Ophelia} and \mention{the Smiths}. \textit{CGEL} describes the proper-noun case as a secondary use within Noun and must treat the compound case as conversion to another category \citep[423, 521]{huddleston2002}. Within Noun, one secondary use covers both, and the article and the adjective follow from what the word now denotes. A separate primary D needs a conversion rule to reach the same facts; the subcategory analysis predicts them.

If determinative-headed and genitive expressions require different structural constraints after their independently motivated restrictions are held fixed, the shared-projection proposal in §\ref{sec:det-uniform} loses its advantage. That would favour separate phrase types. Retaining a separate primary D requires the further case that the category boundary captures the recurring differences better than a determinative subcategory within Noun.

\section{Coordinate or nested subcategories}\label{sec:inheritance}

Including determinatives within Noun leaves a further question: where within Noun do they belong? In the proposed hierarchy, common noun, proper noun, pronoun, and determinative are four coordinate subcategories. Hudson puts his determiner category inside pronoun, which is itself inside Noun. Both group the relevant determinatives with nouns. They differ in whether pronouns and determinatives form an intermediate category that excludes common and proper nouns.

\textcite[253--254]{hudson2010wordgrammar} makes taxonomic inclusion explicit. His noun category contains common noun, proper noun, and pronoun; what he calls a determiner is a pronoun whose \term{valency} permits the relevant common-noun dependent \citep[9--10]{hudson2004determiners}. The word's category remains constant across independent and dependent uses, which differ in their permitted dependents.

Hudson's criteria centre on licensing a singular count common noun and on mutual exclusion in that use. He sets aside \mention{all}, cardinal numerals, and quantifiers restricted to plural or non-count nouns \citep[9--10]{hudson2004determiners}. Their place in his taxonomy remains unsettled by those criteria. Extending the pronoun analysis to \textit{CGEL}'s fuller inventory requires comparing its modifier selection, grade, and quantitative meanings with the existing pronouns' properties.

The intermediate category would be useful if it supported grammatical generalizations applying to pronouns and determinatives together. Coordinate rank permits unequal similarities among the noun subcategories. The issue is whether their shared properties warrant another level in the lexical hierarchy.

\textcite{reynolds2021} compares 138 word forms through properties recorded as present or absent, such as accepting \mention{almost}. Its unsupervised clustering groups forms by these similarities without being given their category labels. The groups broadly resemble the pronoun and determinative inventories, though the outcome varies with initialization and feature selection.\footnote{Relative to a recovered 232-feature working matrix, the public 155-feature file omits 76 singleton columns and one all-zero column. The accompanying \href{run:matrix-audit.pdf}{\textit{Replication audit of the English determinative--pronoun feature matrix}} documents the remaining coding differences, reproduces the published statistical decomposition, and examines sensitivity. The recovered matrix isn't authenticated as the published input; the public file is preserved unchanged.} But the study contains no common or proper nouns. Distinguishing the two groups therefore establishes neither their taxonomic rank nor the proposed superordinate category.

Adding common nouns, proper nouns, and adjectives would require new sampling, feature selection, and coding across all groups. The examples discussed here aren't representative samples of the open categories.

The grammatical case for nesting includes shared lexical properties (Table~\ref{tab:existing}). Pronouns and determinatives both have closed inventories, and both include gender-sensitive forms and interrogative or relative forms. The contrasts \mention{she}/\mention{it} and \mention{somebody}/\mention{something} are encoded by the lexical head, whereas descriptions and names shape referent construal (§\ref{sec:inflection}). These connections are more specific than their general dependence on context.

The \mention{my}/\mention{mine} and \mention{no}/\mention{none} paradigms supply a further grammatical parallel (§\ref{sec:form-selection}). In ordinary NP uses, both distinguish dependent and independent forms. The parallel has limited lexical reach, but it supplies positive evidence for the proposed intermediate grouping.

Modifier selection gives mixed support. Personal pronouns permit restricted adjectival modification, as in \mention{poor old me} \citep[429--430]{huddleston2002}; compare determinative \mention{the lucky few}. Both groups permit peripheral NP modifiers, subject to lexical selection (§\ref{sec:payne}). The four gradable quantifiers also permit the internal degree-AdvP series discussed in §\ref{sec:adjectival-profile}.

Determination gives weaker support for an exclusive pronoun--determinative grouping. Personal pronouns and primary naming uses of proper nouns both resist free determination \citep[429--430, 517, 519--520]{huddleston2002}. Determinatives allow \mention{the few}, \mention{the two}, and \mention{these three}, with item-specific conditions; \textit{CGEL} gives \mention{these few here} and \mention{the many who did} \citep[415--416]{huddleston2002}. These distributions show how restrictions intersect across the noun subcategories.

I retain coordinate rank because the closer affinities cross the proposed intermediate boundary. Partitives, number transparency, restricted dependents, and degree uses connect determinatives with quantificational common nouns (§\ref{sec:quant-nouns}), while reference and selected form alternations connect them with pronouns. The coordinate analysis records both connections within Noun and preserves the distinct organization of each subcategory: person, case, and reflexivity in much of the pronoun system; quantification, determination, and degree in much of the determinative system.

Both inheritance paths meet the fragment's Noun condition (§\ref{sec:det-uniform}), and nesting remains a coherent alternative. The coordinate preference rests on these intersecting profiles; accepting Noun membership and ordinary Head leaves this further choice open.

\section{Conclusion}\label{sec:conclusion}

I propose including English determinatives within Noun. The strongest positive comparison comes from quantificational common nouns: their selected complements, number transparency, restricted dependents, and degree uses connect with parts of the determinative inventory. The connected adjectival profile of the four gradable quantifiers deserves weight on the same terms. Broader constructional and referential connections, with partial inflectional corroboration, favour the nominal grouping.

The grouping includes restricted members through their integration into the determinative system. Articles participate in its contrasts in definiteness, quantity, and count selection; \mention{no}/\mention{none} retains its form-selection conditions.

The ordinary-Head analysis uses the same chain in bare, partitive, and externally determined expressions: the word heads Nom, and Nom heads NP. The taxonomic and structural choices complement each other: the nominal profile supports grouping determinatives with nouns, and that grouping satisfies the proposed grammar's category condition on lexical Heads in Nom. The grouping also unifies two secondary uses that \textit{CGEL} keeps apart: the set-denoting compound (\mention{a special someone}) and the set-denoting proper noun (\mention{an Ophelia}) fall under one process once both are nouns (§\ref{sec:compounds}).

Attributive degree modification still requires a fuller account of lexical, constructional, and register conditions. Coordinate rank alongside pronoun is a further preference. The comparisons connect determinatives with both pronouns and quantificational common nouns.

\appendix

\section{Earlier accounts and logical alternatives}\label{sec:historical}

Keeping a word in one category across dependent and independent uses leaves its taxonomic position open. As \textcite[232--235]{lyons1968} observes, distribution can be compared at different levels: two expressions may belong together at one level and differ at a more specific level. Table~\ref{tab:rivals} distinguishes the selected authors' proposals, their analytical levels, and further logical alternatives.

\begin{table}[H]
\centering\small
\caption{Selected accounts and logical alternatives. The earlier proposals differ in scope and representational level; they don't all specify a surface taxonomy. The final four rows map further possibilities beyond the developed comparisons.}\label{tab:rivals}
\setlength{\tabcolsep}{4pt}
\begin{tabular}{>{\raggedright\arraybackslash}p{2.3cm}>{\raggedright\arraybackslash}p{6.2cm}>{\raggedright\arraybackslash}p{3.05cm}}
\toprule
Account & Proposed relationship & Level or status \\
\midrule
Palmer (\citeyear{palmer1924}) & Determinative adjectives grouped with pronouns. & Lexical grouping; inclusive Noun unspecified \\
Postal (\citeyear{postal1966}) & Personal pronouns analysed as articles, with deeper noun features. & Intermediate and underlying representations \\
Sommerstein (\citeyear{sommerstein1972}) & Definite article and personal pronouns given NP structure. & Underlying representation \\
Lyons (\citeyear{lyons1968}) & Articles, demonstratives, and personal pronouns linked by definiteness and deixis. & Semantic relationship \\
\textit{CGEL} & Pronoun within Noun; determinative separate. & Lexical taxonomy \\
Hudson (\citeyear{hudson2004determiners}) & His determiner category within pronoun, which is within noun (see §\ref{sec:inheritance}). & Lexical taxonomy \\
Spinillo (\citeyear{spinillo2004reconceptualising}) & Determinatives redistributed; \mention{the}, \mention{a}, and \mention{every} retained as articles. & Lexical recategorization \\
D-noun categorization & Common noun, proper noun, pronoun, and determinative coordinate within Noun. & Proposed taxonomy \\
Logical alternative & Noun, pronoun, and determinative separate. & For comparison \\
Logical alternative & Determinative within Noun; pronoun separate. & For comparison \\
Logical alternative & Pronoun within determinative, which is within Noun. & For comparison \\
Logical alternative & Pronoun within proper noun; determinative within common noun. & For comparison \\
\bottomrule
\end{tabular}
\end{table}

\textcite[24]{palmer1924} proposed placing \enquote{determinative adjectives} with pronouns. He contrasted them with qualifying adjectives, which permit predicative use, comparison, and adverbial modification. Most determinatives, he observed, can \enquote{be used indifferently as pronouns or as modifiers of nouns}. \textcite[279]{lyons1968} adds a semantic connection: articles, demonstratives, and personal pronouns share definiteness and deictic contrasts.

\textcite{postal1966} develops the article--pronoun connection through English reflexives and combinations such as \mention{we men}. His category assignments depend on representational level: article segments occur at intermediate stages, their deepest counterparts are features of nouns, and surface pronouns can receive derivative Noun status. This differs from a proposal about the lexical taxonomy of the surface words.

\textcite[197--203]{sommerstein1972} argues in the opposite direction, giving the definite article and personal pronouns underlying NP structure. His English comparison includes the count restriction on anaphoric \mention{one}, whereas \mention{it} can refer to a quantity of a substance.\footnote{\textcite[797--798]{payne2013anaphoric} analyse anaphoric \mention{one} as a common count noun, distinct from determinative \mention{one} and personal pronoun \mention{one}.}

These predecessors challenge a fundamental article--pronoun separation but don't settle the full modern determinative inventory's surface taxonomy. Its contrasts in grade, modification, complementation, and restricted independent use extend beyond their proposals. Hudson's explicit nesting and Spinillo's restricted article grouping are assessed in §§\ref{sec:inheritance} and~\ref{sec:articles}, respectively.

\section*{Data and analysis materials}

The accompanying supplements are \href{run:matrix-audit.pdf}{\textit{Replication audit of the English determinative--pronoun feature matrix}}, \href{run:corpus-documentation.pdf}{\textit{CGELBank concordance and extraction notes}}, and \href{run:quantifier-controls.pdf}{\textit{Quantificational controls: attestations and provenance}}. The \texttt{analysis/} directory preserves their input files, provenance records, scripts, numerical outputs, and sentence concordance. Its README identifies the files and reproduction procedures.

\clearpage
\printbibliography
\end{document}
